%
% 13-komplex.tex -- Der Körper der komplexenzahlen
%
% (c) 2026 Prof Dr Andreas Müller
%
\subsection{Der Körper der komplexen Zahlen}
In der bisherigen Entwicklung haben wir uns vor allem um einzelne
Zahlen gekümmert.
Der Menge der Zahlen als Ganzes und die darauf definierten 
Rechenoperationen haben wir uns nur indirekt angenommen.

%
% Was ist ein Körper
%
\subsubsection{Was ist ein Körper?}
Die moderne Mathematik hat die möglichen Strukturen von Zahlenmengen und
ihren Operationen im Detail studiert und entsprechend ihren Eigenschaften
klassifiziert.
Die einfachste Art einer Struktur geht von nur einer Operation aus.
Die zugehörige Struktur ist die einer Gruppe.

\begin{definition}[Gruppe]
Eine \emph{Gruppe} ist eine Menge $G$ mit einer zweistelligen Operation
\index{Gruppe}%
\[
\cdot
:
G\times G \to G
:
(g,h) \mapsto gh
\]
mit den folgenden Eigenschaften.
\begin{enumerate}
\item
Die Operation ist \emph{assoziativ}: $(gh)l = g(hl)$ für alle $g,h,l\in G$.
\index{assoziativ}%
\item
Es gibt ein \emph{neutrales Element} $e\in G$, es gilt $ea=a$ für alle $a\in G$.
\index{neutrales Element}%
\index{Element!neutral}%
\item
Zu jedem Element $g\in G$ gibt es ein inverses Element $g^{-1}\in G$
mit $g^{-1}g=e$.
\end{enumerate}
Die Gruppe heisst \emph{kommutativ}
\index{kommutativ!Gruppe}%
\index{Gruppe!kommutativ}%
oder \emph{abelsch},
\index{abelsch}%
wenn die Gruppenoperation kommutativ ist.
In einer abelschen Gruppe wird die Gruppenoperation oft additiv geschrieben
und das neutrale Element mit $0$ bezeichnet.
\end{definition}

\begin{beispiel}
Auf der Menge $\mathbb{N}$ ist die Addition definiert und $0\in\mathbb{N}$ 
\index{N@$\mathbb{N}$}%
ist das neutrale Element bezüglich der Addition.
Trotzdem ist $\mathbb{N}$ keine Gruppe, da kein von 0 verschiedenes Element
ein Inverses hat.
In der Menge $\mathbb{Z}$ gibt es zu jeder Zahl $n\in\mathbb{N}$ das
\index{Z@$\mathbb{Z}$}%
inverse Element $-n\in\mathbb{Z}$, daher ist $\mathbb{Z}$ eine
abelsche Gruppe bezüglich der Addition.
\end{beispiel}

\begin{beispiel}
Die reellen Zahlen $\mathbb{R}$ bilden bezüglich der Addition eine
\index{R@$\mathbb{R}$}%
abelsche Gruppe mit dem neutralen Element $0\in\mathbb{R}$.
Die positiven reellen Zahlen $\mathbb{R}^+$ bilden bezüglich der
Multiplikation eine abelsche Gruppe mit neutralem Element
$1\in\mathbb{R}^+$.
Die Exponentialfunktion
\index{Exponentialfunktion}%
\index{exp@$\operatorname{exp}(z)$}%
\index{ez@$e^z$}%
\[
\exp
\colon
\mathbb{R}\to\mathbb{R}^+
:
x\mapsto e^x
\]
bildet die Addition in $\mathbb{R}$ in die Multiplikation in $\mathbb{R}^+$
ab, es gilt
\begin{align*}
\exp(a+b) &= \exp(a)\exp(b),
&
\exp(0)&=1
&&\text{und}
&
\exp(a)^{-1} &= \exp(-a).
\end{align*}
Eine solche Abbildung heisst ein \emph{Gruppenhomomorphismus}.
\index{Homomorphismus!Gruppen-}%
\index{Gruppenhomomorphismus}%
\end{beispiel}

\begin{beispiel}
Die Menge 
\[
\operatorname{GL}_n(\mathbb{R})
=
\{
A\in M_{n\times n}(\mathbb{R})
\mid
\text{$A$ ist invertierbar}
\}
\]
der invertierbaren $n\times n$-Matrizen ist eine Gruppe bezüglich
der Matrizenmultiplikation mit der Einheitsmatrix $I$ als neutralem
Element.
\index{Einheitsmatrix}%
Da die Matrizenmultiplikation für $n>1$ nicht kommutativ ist,
ist $\operatorname{GL}_n(\mathbb{R})$ im Allgemeinen keine abelsche
Gruppe.
\end{beispiel}

\begin{beispiel}
Die Menge der Matrizen
\[
\operatorname{SL}_n(\mathbb{R})
=
\{
A\in M_{n\times n}(\mathbb{R})
\mid
\det A=1
\}
\]
mit Determinante $1$ ist eine Gruppe, da alle Matrizen mit Determinante
$\ne 0$ invertierbar sind.
\index{Determinante}%
Es ist also 
$
\operatorname{SL}_n(\mathbb{R})
\subset
\operatorname{GL}_n(\mathbb{R})
$.

Sogar die Menge
\[
\operatorname{SL}_n(\mathbb{Z})
=
\{
A\in M_{n\times n}(\mathbb{Z})
\mid
\det A = 1
\}
\]
ist eine Gruppe, denn die Einträge der inversen Matrix haben den gemeinsamen
Nenner $\det A$, während der Zähler die Determinante einer Minormatrix ist.
\index{Minormatrix}%
Wegen $\det A=1$ folgt, dass die Einträge von $A^{-1}$ ebenfalls ganzahlig
sind.
\end{beispiel}

Die vertrauten Zahlenmengen zeichnen sich dadurch aus, dass sie nicht
nur eine Operation tragen, sondern zusätzlich zur Addition auch
eine Multiplikation.

\begin{definition}[Ring]
\index{Ring}%
Ein \emph{Ring} ist eine Menge $R$ mit zwei Operationen,
der Addition $+$ und der Multiplikation $\cdot$, mit den folgenden
Eigeschaften:
\begin{enumerate}
\item
Es gibt ein neutrales Element $0\in R$ der Addition so, dass $R$
bezüglich der Addition eine abelsche Gruppe ist.
\item
Die Multiplikation ist assoziativ und \emph{distributiv}:
\index{distributiv}%
\begin{align*}
a(bc) &= (ab)c
\\
a(b+c)&= ab + ac & (a+b)c &= ac + bc.
\end{align*}
\item
$R$ heisst \emph{Ring mit Eins}, wenn es in $R^*=R\setminus\{0\}$ ein
ein neutrales Element $1\in R^*$ der Multiplikation gibt.
\item
Der Ring $R$ heisst \emph{kommutativ}, wenn die Multiplikation kommutativ ist.
\index{kommutativ}%
\item
Eine \emph{Einheit} in $R$ ist eine Element $u\in R^*$, welches bezüglich
\index{Einheit}%
der Multiplikation ist.
\item
Ein Element $a\in R^*$ heisst ein \emph{Nullteiler}, wenn es ein Element
\index{Nullteiler}%
$b\in R^*$ gibt, sodass $ab=0$ ist.
\end{enumerate}
\end{definition}

Nullteiler sind Elemente in $R^*$, die bezüglich der Multiplikation
nicht invertierbar sind.
Ist nämlich $ab=0$ und wäre $a^{-1}$ ein zu $a$ inverses Element, dann
wäre
\[
0
=
a^{-1}0
=
a^{-1}ab
=
b,
\]
ein Widerspruch zu $b\in R^*$.

\begin{beispiel}
Die Menge $\mathbb{Z}$ der ganzen Zahlen ist ein Ring mit Eins.
Die einzigen Elemente, die bezüglich der Multiplikation invertierbar
sind, sind $\pm 1$.
In $\mathbb{Z}$ gibt es keine Nullteiler.
\end{beispiel}

\begin{beispiel}
Sei $n\in\mathbb{N}$ eine natürliche Zahl und
\[
\mathbb{Z}_n
=
\{
0,\dots, n-1
\}
\]
die Menge der Reste bei Teilung durch $n$.
Die Addition und die Multiplikation von ganzen Zahlen überträgt sich
entsprechende Operationen auf $\mathbb{Z}_n$, sodass $\mathbb{Z}_n$
ebenfalls ein kommutativer Ring mit $1$ ist.
\end{beispiel}

\begin{beispiel}
Die Menge $M_{n\times n}(\mathbb{R})$ der $n\times n$-Matrizen mit
reellen Einträgen ist ein Ring mit Eins bezüglich der Addition und
Multiplikation von Matrizen.
Für $n>1$ ist die Matrizenmultiplikation nicht kommutativ, daher
ist $M_{n\times n}(\mathbb{R})$ ein nicht-kommutativer Ring.

Für $n>1$ hat $M_{n\times n}(\mathbb{R})$ viele Nullteiler,
Zum Beispiel gilt
\[
\begin{pmatrix}
1&0\\0&0
\end{pmatrix}
\begin{pmatrix}
0&0\\0&1
\end{pmatrix}
=
\begin{pmatrix}
0&0\\
0&0
\end{pmatrix}
\]
die beiden Matrizen links sind also Nullteiler und können nicht
invertiert werden.
\end{beispiel}

Die Mengen $\mathbb{R}$ und $\mathbb{F}_p$ zeichen sich dadurch
aus, dass nicht nur jedes Element ein additives Inverses hat, sondern
auch jedes von $0$ verschiedene Element ein multiplikatives Inverses hat.
Eine solche Menge heisst ein Körper.

\begin{definition}[Körper]
\index{Körper}%
Ein Körper $K$ ist ein Ring mit Eins derart, dass $K^*$ eine abelsche
Gruppe bezüglich der Multiplikation ist.
\end{definition}

\begin{beispiel}
Die Menge $\mathbb{Q}$ der rationalen Zahlen ist ein Körper.
\index{Q@$\mathbb{Q}$}%
\index{rationale Zahlen}%
\end{beispiel}

\begin{beispiel}
Die Menge $\mathbb{F}_p$ ist für jede Primzahl ein Körper.
\end{beispiel}

%
% Körperaxiome für R
%
\subsubsection{Körperaxiome für $\mathbb{R}$}
Die reellen Zahlen wurden als Grenzwerte von Cauchy-Folgen von
rationalen Zahlen konstruiert.
Sind $(a_n)_{n\in\mathbb{N}}$ und $(b_n)_{n\in\mathbb{N}}$ zwei
Cauchy-Folgen von rationalen Zahlen mit Grenzwerten $a$ bzw.~$b\in\mathbb{R}$.
In der Analysis lernt man, dass
\begin{align*}
\lim_{n\to\infty} (a_n+b_n)
&=
\lim_{n\to\infty} a_n + \lim_{n\to\infty} b_n
=
a+b
\\
\lim_{n\to\infty} (a_nb_n)
&=
\Bigl(\lim_{n\to\infty} a_n\Bigr) \cdot \Bigl(\lim_{n\to\infty} b_n\Bigr)
=
ab
\intertext{und falls $b_n\ne 0$ und $b\ne 0$ gilt auch}
\lim_{n\to\infty} \frac{1}{b_n}
&=
\frac{1}{\displaystyle\lim_{n\to\infty}b_n}
=
\frac{1}{b}
\\
\lim_{n\to\infty}\frac{a_n}{b_n}
&=
\frac{\displaystyle\lim_{n\to\infty} a_n}{\displaystyle\lim_{n\to\infty}b_n}
=
\frac{a}{b}.
\end{align*}
Diese Gleichungen zeigen, dass die Grenzwertbildung, durch die $\mathbb{R}$
definiert wird, mit den arithmetischen Operationen kompatibel sind.
In der Analysis spricht man davon, dass die arithmetischen Operationen
stetig sind.
Insbesondere folgt, dass $\mathbb{R}$ bezüglich der Addition eine
abelsche Gruppe ist und $\mathbb{R}^*$ bezüglich der Multiplikation
eine multiplikative Gruppe.
$\mathbb{R}$ ist also ein Körper, der die rationalen Zahlen als
echten Teilkörper enthält.

%
% Körperaxiome für eine algebraische Erweiterung
%
\subsubsection{Körperaxiome für eine algebraische Erweiterung}
Sei $K$ ein Körper und $m(X)\in K[X]$ ein irreduzibles 
Polynom, welches 
\[
m(X)
=
X^n + m_{n-1}X^{n-1} + \dots + m_1X + m_0
\]
geschrieben werden soll.
Der Körper $K$ kann durch ein Element $\alpha$, welches die
Identität
\begin{equation}
m(\alpha)
=
\alpha^n
+
m_{n-1}\alpha^{n_1}
+
\dots
+
m_1\alpha
+
m_0
\label{buch:komplex:zahlen:eqn:alphaidentitaet}
\end{equation}
erfüllt, erweitert werden.
Dazu werden die Rechenregeln von Polynomen in $\alpha$ mit Koeffizienten
aus $K$ mit der Identität~\eqref{buch:komplex:zahlen:eqn:alphaidentitaet}
reduziert.
Wie früher gezeigt wurde, entsteht so die Menge
\[
K(\alpha)
=
\{
a_0+a_1\alpha + \dots a_{n-1}\alpha^{n-1}
\mid
a_0,\dots,a_{n-1}\in K
\},
\]
die ein $n$-dimensionalers Vektorraum über dem Körper $K$.
Die bereits früher konstruierte Definition der Multiplikation
macht aus $K(\alpha)$ eine Körper, in dem $\alpha$ eine Nullstelle
des Polynoms $m$ ist.

%
% Der Körper C = R(i)
%
\subsubsection{Die Körper $\mathbb{C}=\mathbb{R}(i)$ und $\mathbb{Q}(i)$}
Sei $m(X) = X^2 + 1 \in\mathbb{R}[X]$.
Dieses Polynom ist irreduzibel und definiert daher eine algebraische
Erweiterung von $\mathbb{R}$, die mit $\mathbb{R}(i)$ bezeichnet wird.
Das Element $i$ ist eine Nullstelle von $m$, es gilt also $i^2+1=0$.
Als algebraische Erweiterung von $\mathbb{R}$ ist $\mathbb{C}=\mathbb{R}(i)$ ein
Körper, er heisst der \emph{Körper der komplexen Zahlen}.
\index{komplexe Zahlen}%
\index{C@$\mathbb{C}$}%
Das Element $i$ heisst die \emph{imaginäre Einheit}.
\index{imaginäre Einheit}%
\index{i@$i$}%

Die Menge $\mathbb{R}(i)$ besteht aus den Zahlen
\[
\mathbb{R}(i)
=
\{
a+bi
\mid
a,b\in\mathbb{R}
\}
\]
mit den Ringoperationen
\begin{align*}
\text{Addition:}&&
(a+bi) \pm (c+di)
&=
(a\pm c) + (b\pm d)i
\\
\text{Multiplikation:}&&
(a+bi)(c+di)
&=
ac-bd + (ad+bd)i.
\end{align*}
Für das zu $c+di$ inverse Element muss man wie früher gezeigt
mit $c-di$ erweitern und erhält
\begin{align*}
\frac{1}{c+di}
&=
\frac{1}{c+di}\cdot\frac{c-di}{c-di}
=
\frac{c-di}{c^2+d^2}
=
\frac{c}{c^2+d^2}
+
\frac{-d}{c^2+d^2}i
\\
\text{und}
\qquad
\frac{a+bi}{c+di}
&=
\frac{ac+bd + (-ad+bc)i}{c^2+d^2}
+
\frac{ac+bd}{c^2+d^2}
+
\frac{-ad+bc}{c^2+d^2}i
.
\end{align*}
Die Operationen zur Berechnung des inversen Elementes sind 
für sich genommen nützlich, wir geben ihnen in der folgenden
Definition einen Namen.

\begin{definition}[konjugiert komplexe Zahl]
Für $z=a+bi\in\mathbb{C}$ heisst
$\bar z = a-bi$ die zu $z$ \emph{konjugiert komplexe} Zahl.
\index{konjugiert komplex}%
\end{definition}

Das Polynom $m(X)=X^2+1\in\mathbb{Q}[X]$ ist natürlich auch als
Polynom mit rationalen Koeffizienten irreduzibel.
Daher lässt sich auch $\mathbb{Q}$ durch die imaginäre Einheit
$i$ algebraisch zum Körper $\mathbb{Q}(i)$ erweitern.
Im Gegensatz zu $\mathbb{C}$ lässt sich in $\mathbb{Q}(i)$
exakt rechnen, es ist daher für Computeralgebrasysteme 
unerlässlich.

Die Gleichung $x^2+2=0$ hat in $\mathbb{C}$ die Lösungen
\[
x=
\pm
i
\!\sqrt{2}
\in\mathbb{C},
\]
die allerdings nicht in $\mathbb{Q}(i)$ ist.
Daher muss der Körper $\mathbb{Q}(i)$ erst um $\!\sqrt{2}$ 
algebraisch zu $\mathbb{Q}(i,\!\sqrt{2})$ erweitert werden.
In diesem Körper findet man die beiden Lösungen der Gleichung $x^2-2=0$.

Man beachte, dass sich die Zahlen $i$ und $-i$ mit algebraischen
Mitteln nicht unterscheiden lassen.
Für Zahlen mit reeller Quadratwurzel kann man dank der Ordnungsrelation
in $\mathbb{R}$ die postive von der negativen Quadratwurzel unterschieden.
Für komplexe Zahlen gibt es keine Vergleichsrelation.

